\section{An explicit Green estimate on the certificate's energy scale}
\label{sec:green}

The exterior estimate needed by the certificate concerns the energy norm,
including at an index-one saddle.  Decay of finite-energy critical fields
under stability of the homogeneous lattice is already established in
\cite[Theorem~2.3]{Ehrlacher2016}; that regularity theorem does not require
strong stability of the critical point.  Higher-order far-field conditions
and their approximation properties have also been developed for crystalline
defects \cite{Braun2025,Wei2026}.  Our purpose here is quantitative: we give
an elementary bound with explicit constants, connect it to the same
projections used by the certificate, and retain uniformity under the
distributed coupling perturbations considered in this article.

Identify integer coordinates with physical vertices by
$x\mapsto(6+x_1+\tfrac12,x_2+\tfrac12)$, and use canonical bonds
$(x,i)=(x,x+e_i)$, where $x\in\Z^2$ and $i\in\{1,2\}$.
The point $x$ is the \emph{anchor} of the bond.  Write
$D_i\theta(x)=\theta(x+e_i)-\theta(x)$ and
$\mathcal H=\ell^2(\Z^2\times\{1,2\})$.  Each unoriented bond is counted
once.  Let $\mathsf P:\mathcal H\to D(\X)$ be the orthogonal gradient
projection.  The range is closed because $D:\X\to\mathcal H$ is an
isometry onto its image.

\begin{lemma}[A numerical kernel bound]
\label{lem:green-kernel}
The convolution kernel of $\mathsf P$ satisfies
\begin{equation}
 |\mathsf P_{ij}(z)|\le\frac{8}{(1+\nrm{z}_\infty)^2},
 \qquad z\in\Z^2,\quad i,j\in\{1,2\}.
 \label{eq:green-kernel}
\end{equation}
For each source orientation $j$ and each integer $\rho\ge1$,
\begin{equation}
 \left(\sum_{\nrm{z}_\infty\ge\rho}\sum_{i=1}^2
       |\mathsf P_{ij}(z)|^2\right)^{1/2}
 \le\frac{16\sqrt2}{\rho}.
 \label{eq:green-column}
\end{equation}
Moreover, the norm of evaluation on any bond is exactly
\begin{equation}
 |D_e v|\le 2^{-1/2}\nrm{v}_\X,
 \qquad \nrm{D_e}_{\X^*}=2^{-1/2}.
 \label{eq:edge-evaluation}
\end{equation}
\end{lemma}

\begin{proof}
For $d_i(k)=\exp(\mathrm i k_i)-1$, the Fourier symbol is
$d_i\overline{d_j}/(|d_1|^2+|d_2|^2)$ away from the origin.
Integrating one coordinate exactly gives
$|\mathsf P_{ij}(z)|\le\pi/(2r^2)<2/r^2$ for
$r=\nrm{z}_\infty\ge1$, with the shifts of the mixed coefficients
included.  Appendix~\ref{app:green-kernel} supplies that calculation.
Since $r+1\le2r$, it proves \eqref{eq:green-kernel}; at $z=0$ the bounded
Fourier multiplier gives $|\mathsf P_{ij}(0)|\le1$.
There are $8r$ anchors on the shell of radius $r\ge1$.  Therefore
\[
 \sum_{\nrm{z}_\infty\ge\rho,i}|\mathsf P_{ij}(z)|^2
 \le1024\sum_{r\ge\rho}\frac{r}{(1+r)^4}
 \le1024\int_\rho^\infty t^{-3}\,dt
 =\frac{512}{\rho^2}.
\]
This proves \eqref{eq:green-column}, including both target orientations.
Finally $\mathsf P_{11}(0)=\mathsf P_{22}(0)=1/2$: the two diagonal
symbols sum to one almost everywhere and are exchanged by coordinate
symmetry.  The squared norm of a coordinate functional restricted to
$D(\X)$ is $\mathsf P_{ii}(0)$.  This proves
\eqref{eq:edge-evaluation} by translation invariance.
\end{proof}

\subsection{A stationary source and a dyadic tail inequality}

The following statement also permits noncompact perturbations of the
couplings.  Set
\[
 E_b(\theta)=\sum_e(1+b_e)(1-\cos D_e\theta),
 \qquad b\in\ell^1(\mathcal E),
\]
so that $E_0=\E$.  For an integer anchor radius $R\ge1$, define
\begin{equation}
 T_\theta(R)=
 \left(\sum_{\nrm{x}_\infty\ge R}\sum_{i=1}^2
          |D_i\theta(x)|^2\right)^{1/2}.
 \label{eq:anchor-tail}
\end{equation}

\begin{theorem}[An effective $R^{-1}$ tail]
\label{thm:green}
Let $DE_b(\theta)=0$ in $\X^*$, $\nrm{\theta}_\X\le M$, and suppose
that an integer $R_*\ge1$ and a number $\eta\ge0$ satisfy
\[
 |D_i\theta(x)|\le\eta
 \quad\text{whenever }\nrm{x}_\infty\ge R_*.
\]
Put
\begin{align}
 q_e&=D_e\theta-(1+b_e)\sin(D_e\theta),
 &Q&=\nrm{q}_{\ell^1},\notag\\
 a&=\frac{\eta^2}{6}+\nrm{b}_{\ell^\infty},
 &R_n&=2^nR_*.
 \label{eq:green-source}
\end{align}
If $2a<1$, then
\begin{equation}
 Q\le\frac{M^3}{6}+\nrm{b}_{\ell^1},\qquad
 T_\theta(R_n)\le\frac{A}{R_n}\quad(n\ge0),
 \label{eq:green-tail}
\end{equation}
where
\begin{equation}
 A=\max\left\{R_*T_\theta(R_*),
       \frac{32\sqrt2 Q}{1-2a}\right\}.
 \label{eq:green-A}
\end{equation}
For an arbitrary real radius $R\ge R_*$, the bound
$T_\theta(R)\le2A/R$ holds.
\end{theorem}

\begin{proof}
The global inequality $|t-\sin t|\le|t|^3/6$ and
$\ell^2\subset\ell^3$ give
$\nrm{D\theta-\sin(D\theta)}_{\ell^1}\le M^3/6$.
Also $\nrm{b\sin(D\theta)}_{\ell^1}\le\nrm{b}_{\ell^1}$.
Thus $q\in\ell^1\cap\mathcal H$ and the asserted bound for $Q$ holds.
Criticality states that $(1+b)\sin(D\theta)$ is orthogonal to every
gradient.  Since $D\theta\in D(\X)$,
\begin{equation}
 D\theta=\mathsf Pq.
 \label{eq:green-representation}
\end{equation}
This uses the Riesz formulation in $\X$; it assumes no bounded inverse
of the vertex Laplacian on $\ell^2(\Z^2)$.

The role of force balance in dislocation far-field estimates is examined
in the recent preprint \cite{BraunHudson2026}. Here
\eqref{eq:green-representation} follows directly from stationarity and
the closed gradient projection. If it is rewritten as the vertex equation
$D^*D\theta=D^*q$, then $q\in\ell^1$ implies $D^*q\in\ell^1$ and
$\sum_x(D^*q)(x)=0$: absolute convergence permits pairing the two
opposite endpoint contributions of every bond. Thus this cancellation
is proved from the source representation, rather than inferred from
far-field decay.

For an even integer $R\ge2R_*$, split $q=q_{\rm in}+q_{\rm out}$ at
the anchor radius $R/2$, assigning anchors of radius $R/2$ to the exterior.
For a target anchor of radius at least $R$ and a source anchor of radius
less than $R/2$, their separation is at least $R/2$.
Minkowski's inequality and \eqref{eq:green-column} yield
\[
 \nrm{\mathbf1_{\{\nrm{x}_\infty\ge R\}}
              \mathsf Pq_{\rm in}}_{\mathcal H}
 \le\frac{32\sqrt2}{R}\nrm{q_{\rm in}}_{\ell^1}
 \le\frac{32\sqrt2 Q}{R}.
\]
On the exterior source set,
\[
 |q_e|\le\frac{|D_e\theta|^3}{6}
          +|b_e|\,|\sin D_e\theta|
 \le a|D_e\theta|.
\]
The contraction property of $\mathsf P$ therefore gives
\begin{equation}
 T_\theta(R)\le\frac{32\sqrt2 Q}{R}+aT_\theta(R/2).
 \label{eq:green-recurrence}
\end{equation}
For $A_n=R_nT_\theta(R_n)$ this is
$A_n\le32\sqrt2 Q+2aA_{n-1}$.  The interval $[0,A]$, with $A$ as
in \eqref{eq:green-A}, is invariant under this scalar inequality.
Induction proves \eqref{eq:green-tail}.  More explicitly,
\[
 A_n\le(2a)^nA_0+
       \frac{32\sqrt2 Q\,[1-(2a)^n]}{1-2a}.
\]
If $R_n\le R<R_{n+1}$, monotonicity gives
$T_\theta(R)\le A/R_n\le2A/R$.
\end{proof}

\begin{remark}
No support or moment condition on $b$ is used.  The exterior perturbation
is $b_e\sin(D_e\theta)$, which is controlled \emph{multiplicatively} by
$D_e\theta$.  Replacing it by an independent additive source would lose
this property and would require additional tail information.
\end{remark}

\subsection{The same compact projections and the upper path}

Let $\X_S$ contain the configurations with a representative zero outside
the integer vertex box $\Lambda_S=[-S,S]^2$, and let
$\Pi_S:\X\to\X_S$ be the energy-orthogonal projection.
The distinction between anchor and endpoint radii is needed when these
projections are compared with stored exterior bounds.

\begin{lemma}[A cutoff on the common energy scale]
\label{lem:green-projection}
For every $\theta\in\X$ and integer $R\ge1$,
\begin{equation}
 \nrm{(I-\Pi_{2R})\theta}_\X
 \le\sqrt{227}
 \left(\sum_{\min(\nrm{x}_\infty,\nrm{y}_\infty)\ge R}
                   |D_{\{x,y\}}\theta|^2\right)^{1/2}
 \le\sqrt{227}\,T_\theta(R).
 \label{eq:green-projection}
\end{equation}
In particular, under Theorem~\ref{thm:green}, the left side at $R=R_n$
is at most $\sqrt{227}A/R_n<16A/R_n$.
\end{lemma}

\begin{proof}
On $\mathcal A_R=\{x:R\le\nrm{x}_\infty\le2R\}$, the discrete
Poincar\'e inequality has constant $(225/4)R^2$
(Appendix~\ref{app:green-annulus}).  Let $c$ be the annular mean and let
$\zeta_R$ equal one up to radius $R$, zero from radius $2R$ onwards,
and interpolate linearly between them.  The compact representative
$\zeta_R(\theta-c)$ gives an admissible approximation in $\X_{2R}$.
Using the exact product identity
$D(fg)=\overline f\,Dg+\overline g\,Df$, where bars denote endpoint
averages, and
$2|\overline f|^2\le|f_x|^2+|f_y|^2$, its squared error is bounded by
\[
 2\sum_{\min(\nrm{x}_\infty,\nrm{y}_\infty)\ge R}|D_e\theta|^2
 +\frac4{R^2}\sum_{\mathcal A_R}|\theta-c|^2
 \le227\sum_{\min(\nrm{x}_\infty,\nrm{y}_\infty)\ge R}|D_e\theta|^2.
\]
An edge below this minimum-endpoint radius has zero error gradient:
integer neighboring radii differ by at most one and the cutoff is still
one at radius $R$.  Best approximation proves the first inequality.
Every edge in that tail has its canonical anchor of radius at least $R$,
which proves the second.
\end{proof}

For the homogeneous certificate, the known upper path is polygonal.
Its only implicit vertices are the exact minimum and the selected exact
saddle.  Each remaining vertex, except the zero endpoint, is $\phi_{12}$ plus a
stored compact correction. Here
\[
 \phi_{12}(x)=\arg(6+x_1+\tfrac12+\mathrm i(x_2+\tfrac12))
 -\arg(-6+x_1+\tfrac12+\mathrm i(x_2+\tfrac12)).
\]
The argument branches cancel their far-field jumps outside the segment
between the two poles.  Once $R\ge16$ and the projection box contains those supports, their
projection errors obey the direct dipolar bound $480/R$, proved in
Appendix~\ref{app:green-inputs}.
Consequently, if $A$ bounds both critical vertices in
Theorem~\ref{thm:green}, linearity and convexity of the norm give
\begin{equation}
 \sup_t\nrm{(I-\Pi_{2R_n})\Gamma_0(t)}_\X
 \le\frac{C_\Gamma}{R_n},\qquad
 C_\Gamma=\max\{16A,480\}.
 \label{eq:green-path}
\end{equation}
No stationary equation is assigned to a segment between vertices.

For finite-box barriers, restrict the paths in \eqref{eq:barrier} to
$\X_S$ and the endpoints to $\K\cap\X_S$, retaining all bonds crossing
the vertex-box boundary. Denote the resulting barrier by $B_S(u)$ and
write $B_\infty=B_0$, $C_\infty(u)=\E(u)+B_0(u)$.
For a compact origin $u\in\X_{2R_n}$, suppose that its segment to $m_0$
lies below the certified level $h$, and that the separator gives
$C_\infty(u)=h$.  If $M_\Gamma$ bounds the norm on that segment and on
$\Gamma_0$, then
\begin{equation}
 0\le B_{2R_n}(u)-B_\infty(u)
 \le\frac{M_\Gamma C_\Gamma}{R_n}.
 \label{eq:green-barrier-localization}
\end{equation}
Indeed, concatenate the segment and $\Gamma_0$ and apply $\Pi_{2R_n}$.
The original upper path ends at zero: its terminal check gives
\texttt{compact\_integer\_lift\_max}=0.
Thus the origin and that endpoint are unchanged. Since
$\nrm{\nabla_\X\E(z)}_\X\le\nrm{z}_\X$ and $\Pi_{2R_n}$ is a
contraction, the energy change of each projected point is at most
$M_\Gamma$ times its projection error.  This yields the upper bound;
the inclusion of finite-box paths in the infinite path class gives the
lower bound.  Thus the comparison uses the same strong energy topology
and keeps all boundary bonds of the finite box.

\subsection{Numerical consequences and uniform perturbed tails}

The payload inequalities in Appendix~\ref{app:green-inputs} permit
$M=M_\Gamma=6007$ for the certified path and for compact origins satisfying
$\nrm u_\X\le6007$, with exterior minimum-endpoint radius $4096$, and
$\eta=9/500$.  The original upper path uses $m_0$ and $s_+$, so anchor
radius $R_*=4097$ is sufficient.  Replacing $32\sqrt2$ by $46$ and
$\sqrt{227}$ by $16$ gives the rational constants
\begin{equation}
 A\le\frac{1246352073472250000}{749919},\qquad
 M_\Gamma C_\Gamma
 \le\frac{119789390485564892000000}{749919}.
 \label{eq:green-data-constants}
\end{equation}
Exact rational arithmetic gives the following budgets for the allocation
$M_\Gamma C_\Gamma/R_n\le\varepsilon/4$:
\begin{center}
\begin{tabular}{cclc}
\hline
$\varepsilon$ & $n$ & $R_n=4097\,2^n$ & error upper bound\\
\hline
$10^{-2}$ & $54$ & $73804990693347688448$ & $<2.165\,10^{-3}$\\
$10^{-6}$ & $68$ & $1209220967519808527532032$ & $<1.321\,10^{-7}$\\
\hline
\end{tabular}
\end{center}
The actual projection box has radius $2R_n$.
These are arithmetic bounds, not newly computed barrier intervals.
They replace the thousands-of-digits radii from the weak hole-filling
estimate by $20$- and $25$-digit values of $R_n$, but still cannot be
materialized.  A full cubical covering would also retain its exponential
dependence on the finite dimension.  Useful radii require sharper
data-dependent bounds for $Q$ and $T_\theta(R_*)$ and the retained
low-dimensional reaction reduction.

The same estimate is uniform over the nearby critical families of the
robustness theorem.  For all three critical states, including the reflected
saddle, use the common values
\begin{equation}
 \nrm{b}_{\ell^1}\le\frac1{500},\quad M=6008,\quad
 \eta=\frac{41}{1000},\quad R_*=4098.
 \label{eq:green-robust-data}
\end{equation}
Here a displacement of at most $1/32$ in $\X$ changes every bond
increment by at most $1/(32\sqrt2)$, so
$0.018+1/(32\sqrt2)<0.041$.  The extra radial unit is deliberate:
the physical reflection $x_1\mapsto12-x_1$ sends integer indices
$i\mapsto-i-1$, and a reoriented horizontal bond anchor to $-i-2$.
An anchor radius of $4098$ therefore guarantees reflected minimum
endpoint radius $4096$.
With \eqref{eq:green-robust-data},
\[
 a\le\frac{13681}{6000000},\qquad
 Q\le\frac{54216288128003}{1500},\qquad
 A\le\frac{4987898507776276000}{2986319}<1.7\,10^{12}.
\]
In particular, all these exact critical states satisfy, on the common
dyadic sequence,
\begin{equation}
 T_\theta(R_n)\le\frac{1.7\,10^{12}}{R_n},\qquad
 \nrm{(I-\Pi_{2R_n})\theta}_\X
 \le\frac{2.72\,10^{13}}{R_n}.
 \label{eq:green-robust-tail}
\end{equation}
This last conclusion allows noncompact $\ell^1$ coupling perturbations.
It is a uniform mathematical tail statement; an implementation accepting
arbitrary infinite inputs still requires an effective representation of
their coefficients and the relevant membership bounds.
